\documentclass{article}
\usepackage[T1]{fontenc}
\usepackage{lmodern}
\usepackage{graphicx}
\usepackage{amsmath,amssymb}
\usepackage[a4paper,margin=2cm]{geometry}
\usepackage[absolute,overlay]{textpos}
\setlength{\TPHorizModule}{1mm}
\setlength{\TPVertModule}{1mm}
\usepackage[indonesian]{babel}
\usepackage{hyperref}
\setlength{\emergencystretch}{2em}
% unit: Halaman 20

\begin{document}

% === Halaman 20 (llm) ===

\textbf{Contoh 5:} Kita akan mendekripsikan cipherteks dari Contoh 4 dengan menggunakan kunci privat $\{2, 3, 6, 13, 27, 52\}$. Di sini, $n = 31$ dan $m = 105$. Nilai $31^{-1} \pmod{105}$ diperoleh sbb:

\[ n \cdot n^{-1} \equiv 1 \pmod{m} \rightarrow 31 \cdot n^{-1} \equiv 1 \pmod{105} \rightarrow n^{-1} = (1 + 105k)/31 \]
coba $k = 0, 1, 2, \dots$, diperoleh $n^{-1} = 61$

Cipherteks dari Contoh 4 adalah $174, 280, 333$. Hasil dekripsi:

\begin{align*}
174 \cdot 61 \pmod{105} = 9 &= \color{red}{0} \cdot 2 + \color{red}{1} \cdot 3 + \color{red}{1} \cdot 6 + \color{red}{0} \cdot 13 + \color{red}{0} \cdot 27 + \color{red}{0} \cdot 52 \rightarrow \color{red}{011000} \\
280 \cdot 61 \pmod{105} = 70 &= \color{red}{1} \cdot 2 + \color{red}{1} \cdot 3 + \color{red}{0} \cdot 6 + \color{red}{1} \cdot 13 + \color{red}{0} \cdot 27 + \color{red}{1} \cdot 52 \rightarrow \color{red}{110101} \\
333 \cdot 61 \pmod{105} = 48 &= \color{red}{1} \cdot 2 + \color{red}{0} \cdot 3 + \color{red}{1} \cdot 6 + \color{red}{1} \cdot 13 + \color{red}{1} \cdot 27 + \color{red}{0} \cdot 52 \rightarrow \color{red}{101110} 
\end{align*}

Jadi, plainteks yang dihasilkan kembali adalah:
\[ \color{red}{011000110101101110} \]

\end{document}
